\documentclass[12pt,a4paper]{article}

% --- PACOTES BÁSICOS ---
\usepackage[utf8]{inputenc}
\usepackage[T1]{fontenc}
\usepackage[brazilian]{babel}
\usepackage{amsmath, amsfonts, amssymb}
\usepackage{graphicx}
\usepackage{geometry}
\usepackage{pdflscape}
\usepackage{booktabs}
\usepackage{pgffor}
\usepackage{xifthen} % <--- melhora testes condicionais
\usepackage[alf]{abntex2cite}

\geometry{a4paper, left=3cm, top=3cm, right=2cm, bottom=2cm}

\newcommand{\enunciado}[1]{\noindent \textit{\large #1}}

% --- DOCUMENTO ---
\begin{document}

% --- CAPA ---
\begin{titlepage}
    \begin{center}
        \textbf{Centro Federal de Educação Tecnológica de Minas Gerais (CEFET-MG)}\\[0.2cm]
        \textbf{Programa de Pós-Graduação em Modelagem Matemática e Computacional}\\[3cm]

        \large{Disciplina: Princípios de Modelagem Matemática}\\[3cm]

        \LARGE\textbf{Principles of Mathematical Modeling (2nd Edition):\\
            Solução dos Problemas 3.1 a 3.36}\\[3cm]

        \large{
        Alexandre Oliveira, Ariel, Ectori, Erasmo, Frederico,\\
        Glauber Silveira, Israel, Jeancarlo Campos Leão, Júlia,\\
        Leonardo V. Guimarães, Marcus Santos, Renan Machado,\\
        Rodrigo Favato, Sergio.
        }\\[3cm]

        \large{Belo Horizonte -- MG -- Brazil}\\
        \large{2026}
    \end{center}
\end{titlepage}

% --- CONTEÚDO ---
\foreach \i in {1,...,36} {

    \section*{Problema 3.\i}

    % Enunciado
    \IfFileExists{enunciado/e\i.tex}{
        \noindent \large \textit{\input{enunciado/e\i.tex}}
    }{
        \textbf{Enunciado não encontrado para o Problema 3.\i}
    }

    \vspace{1cm}

    % Resposta
    \IfFileExists{resposta/r\i.tex}{
        \input{resposta/r\i.tex}
    }{
        \textbf{Resposta não encontrada para o Problema 3.\i}
    }

    \newpage
}


\bibliographystyle{abntex2-alf}
\bibliography{references}

\end{document}

